% ============================================================
% Paper 93 (EN): Gromov's Filling Radius and Systolic Geometry
% Author: Michael Fuhrmann
% Project: Specialist Mathematics Project
% Build: 169
% Date: 2026-03-12
% ============================================================
\documentclass[12pt,a4paper]{amsart}

\usepackage{amsmath,amssymb,amsthm,mathtools}
\usepackage{enumitem}
\usepackage{booktabs}
\usepackage{array}
\usepackage[hidelinks]{hyperref}
\usepackage[utf8]{inputenc}
\usepackage[T1]{fontenc}
\usepackage{geometry}
\geometry{margin=2.5cm}

% ---- Theorem Environments -----------------------------------
\newtheorem{theorem}{Theorem}[section]
\newtheorem{lemma}[theorem]{Lemma}
\newtheorem{corollary}[theorem]{Corollary}
\newtheorem{proposition}[theorem]{Proposition}
\theoremstyle{definition}
\newtheorem{definition}[theorem]{Definition}
\newtheorem{example}[theorem]{Example}
\theoremstyle{remark}
\newtheorem{remark}[theorem]{Remark}
\newtheorem{conjecture}[theorem]{Conjecture}

% ---- Custom Operators ---------------------------------------
\DeclareMathOperator{\sys}{sys}
\DeclareMathOperator{\FillRad}{FillRad}
\DeclareMathOperator{\vol}{vol}
\DeclareMathOperator{\diam}{diam}
\DeclareMathOperator{\inj}{inj}
\newcommand{\RR}{\mathbb{R}}
\newcommand{\ZZ}{\mathbb{Z}}
\newcommand{\RP}{\mathbb{RP}}
\newcommand{\TT}{\mathbb{T}}

% ---- Abstract -----------------------------------------------
\begin{abstract}
Systolic geometry studies the relationship between the shortest
non-contractible loop (the systole) and the volume of a Riemannian manifold.
Gromov's fundamental inequality $\sys(M)^n \le C_n \vol(M)$ holds for all
\emph{essential} $n$-manifolds.  The key intermediate quantity is the
\emph{filling radius} $\FillRad(M)$, and Gromov proved the bound
$\FillRad(M) \le C_n \sys(M)$.  We treat the classical inequalities of
Pu (for $\RP^2$) and Loewner (for $\TT^2$), the Gromov compactness theorem
for metric spaces, and the role of Morse theory for the length functional.
Sharp systolic constants are known only in special cases; many remain open.
\end{abstract}

\title{Gromov's Filling Radius and Systolic Geometry}

\author{Michael Fuhrmann}
\address{Specialist Mathematics Project, Build 169}
\email{specialist-maths@project.local}
\date{2026-03-12}

\maketitle
\tableofcontents

% =============================================================
\section{Introduction}
% =============================================================

One of the oldest questions in Riemannian geometry is: how short can a
closed curve be forced to be by topological constraints?  Given a closed
Riemannian manifold $(M,g)$, the \emph{systole} $\sys(M,g)$ is the
shortest length among all non-contractible closed curves in $M$.  A natural
question then is how $\sys(M,g)$ relates to geometric data such as
$\vol(M,g)$.

Carl Loewner proved in 1949 (unpublished, attributed in \cite{Pu1952}) that
for any metric $g$ on the 2-torus $\TT^2$,
\[
  \sys(\TT^2, g)^2 \le \frac{2}{\sqrt{3}} \vol(\TT^2, g),
\]
with equality for the flat equilateral torus.  Pu \cite{Pu1952} proved
an analogous inequality for $\RP^2$.

In 1983 Gromov \cite{Gromov1983} proved a sweeping generalisation: for
every \emph{essential} closed $n$-manifold $M$,
\[
  \sys(M)^n \le C_n \vol(M).
\]
The proof introduced a new geometric quantity, the \emph{filling radius}
$\FillRad(M)$, and went via the isoperimetric inequality $\FillRad \le C_n \sys$.

This paper presents a rigorous treatment of these ideas, discussing also
the Sabourau--Rotman improvements, Morse theory for the loop space, and the
many open problems concerning sharp constants.

% =============================================================
\section{The Systole}
% =============================================================

\begin{definition}[Systole]
Let $(M,g)$ be a closed Riemannian manifold with $\pi_1(M) \neq 1$ (or more
generally with non-trivial $\pi_k(M)$ for some $k$).  The
\emph{1-systole} is
\[
  \sys_1(M,g) = \inf\bigl\{ \mathrm{length}(\gamma) : \gamma \text{ is a
  closed non-contractible loop in } M \bigr\}.
\]
More generally, the \emph{$k$-systole} $\sys_k(M,g)$ is the infimum of
volumes of $k$-cycles not bounding in $M$.
\end{definition}

\begin{remark}
In this paper ``systole'' always means $\sys_1$ unless stated otherwise.
The infimum is always achieved (by a geodesic) when $M$ is compact.
\end{remark}

\begin{example}
\begin{enumerate}[label=(\roman*)]
  \item For the flat torus $\TT^2 = \RR^2/\Lambda$, $\sys = \lambda_1(\Lambda)$
        (shortest lattice vector).
  \item For $\RP^n$ with the round metric of circumradius $R$,
        $\sys = \pi R$ (the projective line).
  \item For a surface of genus $g \ge 2$ with hyperbolic metric,
        $\sys \le 2\log(4g-2)$ by a collar lemma argument.
\end{enumerate}
\end{example}

% =============================================================
\section{The Filling Radius}
% =============================================================

Gromov's key insight was to measure how "deep" a manifold is buried inside
any ambient metric space.

\begin{definition}[Filling Radius \cite{Gromov1983}]
Let $M$ be a closed orientable $n$-manifold embedded isometrically in the
Banach space $L^\infty(M)$ via the Kuratowski embedding
$\iota \colon x \mapsto d(x,\cdot)$.  The \emph{filling radius} is
\[
  \FillRad(M) = \inf\bigl\{ r > 0 : \iota_*[M] = 0 \in H_n(U_r(M);\ZZ)
  \bigr\},
\]
where $U_r(M) = \{f \in L^\infty(M) : d(f, \iota(M)) < r\}$ is the
$r$-tubular neighbourhood of $\iota(M)$ in $L^\infty(M)$, and $[M]$ is
the fundamental class.
\end{definition}

\begin{remark}
This definition is independent of the choice of ambient Banach space,
as long as the metric is isometrically embedded.  The Kuratowski embedding
$\iota(x) = d(x,\cdot) \in L^\infty(M)$ is always an isometric embedding.
\end{remark}

\begin{example}
\begin{enumerate}[label=(\roman*)]
  \item For a round sphere $S^n$ of radius $R$: $\FillRad(S^n) = \arccos(-\tfrac{1}{n+1}) \cdot R$
        (by Katz's computation \cite{Katz1983}).
  \item For the flat torus $\TT^n$: $\FillRad(\TT^n) \le c_n \sys(\TT^n)$.
  \item For $\RP^n$ with the round metric: $\FillRad = \tfrac{\pi}{3} R$.
\end{enumerate}
\end{example}

% =============================================================
\section{Essential Manifolds}
% =============================================================

\begin{definition}[Essential Manifold]
A closed $n$-manifold $M$ is \emph{essential} if the classifying map
$f \colon M \to K(\pi_1(M), 1)$ satisfies $f_*[M] \neq 0 \in
H_n(K(\pi_1(M),1);\ZZ)$.
\end{definition}

\begin{example}
\begin{enumerate}[label=(\roman*)]
  \item Surfaces of genus $g \ge 1$ are essential.
  \item $\RP^n$ is essential (the fundamental class is non-zero in
        $H_n(K(\ZZ/2,1);\ZZ/2)$).
  \item $S^n$ is \emph{not} essential (since $\pi_1(S^n) = 1$ for $n \ge 2$).
  \item Products of essential manifolds are essential.
  \item Aspherical manifolds (universal cover contractible) are essential.
\end{enumerate}
\end{example}

Essential manifolds are exactly the class for which the systolic inequality
is non-trivial.  For non-essential manifolds the systole can be made
arbitrarily small by rescaling, while the volume stays bounded.

% =============================================================
\section{Gromov's Systolic Inequality}
% =============================================================

\begin{theorem}[Gromov, 1983 \cite{Gromov1983}]
\label{thm:gromov-systolic}
For every essential closed Riemannian $n$-manifold $(M,g)$,
\[
  \sys_1(M,g)^n \le C_n \cdot \vol(M,g),
\]
where $C_n > 0$ is a constant depending only on $n$.
\end{theorem}

Gromov's proof proceeds via the filling radius in two main steps:

\begin{theorem}[Filling Radius Bound \cite{Gromov1983}]
\label{thm:fillrad-sys}
For every closed Riemannian $n$-manifold $(M,g)$,
\[
  \FillRad(M,g) \le \frac{1}{2} \sys_1(M,g).
\]
More precisely, $\FillRad(M,g) \le \tfrac{n+1}{2(n+2)} \sys_1(M,g)$.
\end{theorem}

\begin{theorem}[Volume Filling Inequality \cite{Gromov1983}]
\label{thm:vol-fillrad}
For every essential closed Riemannian $n$-manifold $(M,g)$,
\[
  \FillRad(M,g)^n \le C_n' \cdot \vol(M,g),
\]
where $C_n' > 0$ depends only on $n$.
\end{theorem}

Theorem~\ref{thm:gromov-systolic} follows immediately by combining
Theorems~\ref{thm:fillrad-sys} and~\ref{thm:vol-fillrad}.

\begin{proof}[Proof of Theorem~\ref{thm:fillrad-sys} sketch]
Let $\gamma$ be a shortest non-contractible loop of length $L = \sys(M)$.
If $\FillRad(M) > L/2$, then the loop $\gamma$ remains non-contractible in
the $L/2$-neighbourhood $U_{L/2}(\iota(M))$ inside $L^\infty(M)$.  By
a nerve argument (using the Vietoris--Rips complex at scale $\varepsilon$
and letting $\varepsilon \to L/2$), one shows that the fundamental class
$[M]$ must already die at radius $\le L/2$, contradiction.
\end{proof}

\begin{proof}[Proof of Theorem~\ref{thm:vol-fillrad} sketch]
The key ingredient is a co-dimension-one isoperimetric inequality in
$L^\infty(M)$.  Gromov uses a \emph{waist inequality}: if the $n$-cycle
$[M]$ is non-trivial, any filling has volume at least proportional to
$\FillRad^n$.  The precise argument involves parametric Morse theory on the
distance function in $L^\infty(M)$.
\end{proof}

% =============================================================
\section{Pu's Inequality for $\RP^2$}
% =============================================================

\begin{theorem}[Pu, 1952 \cite{Pu1952}]
For any Riemannian metric $g$ on $\RP^2$,
\[
  \sys(\RP^2, g)^2 \le \frac{\pi}{2} \vol(\RP^2, g),
\]
with equality if and only if $(\RP^2, g)$ is the standard round metric of
constant curvature $+1$ (normalised appropriately).
\end{theorem}

\begin{proof}[Proof idea]
Pu's proof uses the uniformization theorem to conformally parameterise
$(\RP^2, g)$.  Write $g = e^{2u} g_{\mathrm{round}}$.  The systole is bounded
below by the conformal length of a non-contractible geodesic, and the
Cauchy--Schwarz inequality yields the bound.  The extremal metric is the
round metric by an equality case analysis.
\end{proof}

The constant $\pi/2$ is called the \emph{Pu constant} and equals
$(\text{area of }\RP^2)/(\text{sys}^2)$ for the round metric.

% =============================================================
\section{Loewner's Torus Inequality}
% =============================================================

\begin{theorem}[Loewner, 1949, cf.\ \cite{Pu1952}]
For any Riemannian metric $g$ on the 2-torus $\TT^2$,
\[
  \sys(\TT^2, g)^2 \le \frac{2}{\sqrt{3}} \vol(\TT^2, g),
\]
with equality if and only if $(\TT^2, g)$ is (up to scaling) the flat torus
$\RR^2 / \Lambda$ where $\Lambda$ is the hexagonal lattice.
\end{theorem}

The constant $2/\sqrt{3}$ is the square of the systole of the hexagonal
lattice divided by its fundamental domain area.  The hexagonal lattice
$\Lambda = \ZZ(1,0) + \ZZ(1/2, \sqrt{3}/2)$ achieves the extremum.

\begin{proof}[Proof sketch]
By uniformization, $g = e^{2u} g_\Lambda$ for some flat metric $g_\Lambda$.
The inequality $\sys^2 \le C \vol$ follows from the lattice geometry:
among all rank-2 lattices of given co-volume, the shortest vector is
maximised by the hexagonal lattice (Lagrange's theorem on lattice reduction).
\end{proof}

% =============================================================
\section{Berger's Conjecture and Sharp Constants}
% =============================================================

\begin{definition}[Systolic Ratio]
The \emph{systolic ratio} of a closed $n$-manifold $M$ with metric $g$ is
\[
  \mathrm{SR}(M,g) = \frac{\sys(M,g)^n}{\vol(M,g)}.
\]
The \emph{optimal systolic ratio} is $\mathrm{SR}(M) = \sup_g \mathrm{SR}(M,g)$.
\end{definition}

The sharp constants $C_n$ in Gromov's inequality are generally not known.
The following sharp results are known:

\begin{center}
\begin{tabular}{lll}
\toprule
Manifold $M$ & $\mathrm{SR}(M)$ & Extremal metric \\
\midrule
$\TT^2$ & $2/\sqrt{3}$ & Hexagonal flat torus (Loewner) \\
$\RP^2$ & $\pi/2$ & Round metric (Pu) \\
$K(\ZZ,2) = \TT^3$ & $?$ & Open (Berger) \\
$\RP^3$ & $?$ & Open \\
\bottomrule
\end{tabular}
\end{center}

\begin{conjecture}[Berger's Problem on Sharp Systolic Constants, Open]
Determine the optimal systolic ratio $\mathrm{SR}(M)$ for all closed
Riemannian manifolds $M$.  In particular, identify the extremal metrics
for $\TT^n$ ($n \ge 3$), $\RP^n$ ($n \ge 3$), and all higher-genus surfaces.
\end{conjecture}

For surfaces of genus $g \ge 2$, the systolic ratio satisfies
$\mathrm{SR}(\Sigma_g) \ge C/g$ (Gromov), but the sharp constant is unknown.

% =============================================================
\section{Morse Theory for the Loop Space}
% =============================================================

Morse theory applied to the energy functional $E(\gamma) = \int_0^1 |\dot\gamma|^2 dt$
on the free loop space $\Lambda M$ provides a powerful tool in systolic geometry.

\begin{theorem}[Ljusternik--Schnirelmann for Geodesics]
Every closed Riemannian manifold $M$ with $\pi_1(M) \neq 1$ contains a
\emph{closed geodesic} of length exactly $\sys_1(M)$.  If $M$ is essential,
every Riemannian metric on $M$ has at least $\mathrm{cat}(M)$ non-trivial
homotopy classes of closed curves, each containing a geodesic.
\end{theorem}

\begin{definition}[Lusternik--Schnirelmann Category]
The \emph{L-S category} $\mathrm{cat}(M)$ is the minimum number of
contractible open sets needed to cover $M$.
\end{definition}

\begin{theorem}[Gromov's Waist Inequality \cite{Gromov2003waist}]
Let $M$ be a closed Riemannian $n$-manifold and $f \colon M \to \RR^k$ a
continuous map.  Then there exists a fibre $f^{-1}(t)$ such that
\[
  \sys(f^{-1}(t))^{n-k} \le C_{n,k} \cdot \vol_k(M)^{(n-k)/n}.
\]
\end{theorem}

% =============================================================
\section{Sabourau--Rotman Improvements}
% =============================================================

\begin{theorem}[Rotman, 2006 \cite{Rotman2006}]
Let $M$ be a closed Riemannian $n$-manifold with $\pi_1(M) \neq 1$.
Then $M$ contains a closed geodesic $\gamma$ (a shortest non-contractible
closed geodesic) of length $\ell(\gamma) = \sys(M)$, satisfying
\[
  \ell(\gamma) \le 2 \diam(M).
\]
\end{theorem}

\begin{theorem}[Sabourau, 2004 \cite{Sabourau2004}]
For every aspherical surface $(M^2, g)$ (i.e., genus $g \ge 1$),
\[
  \sys(M)^2 \le \gamma(M) \cdot \vol(M),
\]
where $\gamma(M) > 0$ is a constant depending only on the topology of $M$,
and equality is achieved in the limit by sequences of metrics approaching
certain optimal shapes.
\end{theorem}

\begin{conjecture}[Optimal Filling Ratio, Open]
\label{conj:filling}
The sharp constant in $\FillRad(M)^n \le C \cdot \vol(M)$ (for essential
$M$) is achieved by a specific "round" or "regular" metric.  In particular,
for $\RP^n$ the sharp filling radius constant is $\tfrac{\pi}{3}$ (the value
for the round metric).
\end{conjecture}

% =============================================================
\section{Katz's Results on Filling Radius}
% =============================================================

\begin{theorem}[Katz, 1983 \cite{Katz1983}]
The filling radius of the round $n$-sphere $S^n(R)$ of radius $R$ is
\[
  \FillRad(S^n(R)) = R \arccos\!\left(-\frac{1}{n+1}\right).
\]
For large $n$ this is approximately $R(\pi/2 - 1/\sqrt{n})$.
\end{theorem}

\begin{theorem}[Katz, 1983 \cite{Katz1983}]
For $\RP^n$ with the round metric of radius $R$:
\[
  \FillRad(\RP^n) = \frac{\pi R}{3}.
\]
\end{theorem}

\begin{corollary}
For $\RP^2$ with any metric $g$ with $\vol(\RP^2, g) = \pi$ (the round
volume):
\[
  \FillRad(\RP^2, g) \le \frac{\pi}{3}
\]
with equality for the round metric, consistent with Pu's inequality.
\end{corollary}

% =============================================================
\section{Gromov Compactness and Geometric Convergence}
% =============================================================

A crucial tool underlying all these results is:

\begin{theorem}[Gromov Compactness Theorem \cite{Gromov1981}]
The class of Riemannian $n$-manifolds $(M,g)$ with
\begin{enumerate}[label=(\roman*)]
  \item Ricci curvature $\mathrm{Ric}(M,g) \ge -(n-1)\kappa^2$,
  \item diameter $\diam(M,g) \le D$,
\end{enumerate}
is precompact in the Gromov--Hausdorff topology: every sequence has a
subsequence converging (in the GH sense) to a compact metric space.
\end{theorem}

This theorem is used in systolic geometry to take limits of sequences of
metrics and study the behaviour of the systole and filling radius in
degenerating families.

% =============================================================
\section{Open Problems in Systolic Geometry}
% =============================================================

We collect the major open problems:

\begin{conjecture}[Sharp Systolic Constant for $\TT^n$, $n \ge 3$]
The optimal systolic ratio for the $n$-torus $\TT^n$ equals
\[
  \mathrm{SR}(\TT^n) = \gamma_n^{n/2},
\]
where $\gamma_n$ is the \emph{Hermite constant} (the maximum ratio
$\lambda_1(\Lambda)^n/\mathrm{covol}(\Lambda)$ over all rank-$n$ lattices
$\Lambda$).  This is known for $n \le 8$ and $n = 24$ (from lattice theory),
but open for general $n$.
\end{conjecture}

\begin{conjecture}[Gromov's Optimal Filling Radius Conjecture, Open]
For an essential $n$-manifold $M$, the sharp constant in
$\sys(M)^n \le C \cdot \vol(M)$ is related to the geometry of the
regular simplex in $\RR^n$.  The precise sharp constant in Gromov's
Theorem~\ref{thm:gromov-systolic} is unknown for $n \ge 3$.
\end{conjecture}

\begin{conjecture}[Systolic Freedom, Open]
For simply-connected manifolds of dimension $\ge 4$, the systolic
ratio $\mathrm{SR}(M,g)$ can be made arbitrarily large (Gromov's
``systolic freedom'' phenomenon).  Characterising which manifolds
exhibit systolic freedom remains open.
\end{conjecture}

% =============================================================
\section{Conclusion}
% =============================================================

Systolic geometry, initiated by Loewner, Pu, and Gromov, provides deep
connections between metric geometry, algebraic topology, and the theory of
Banach spaces.  The filling radius---the measure of how deeply a manifold is
embedded in its Kuratowski completion---turns out to be the key intermediate
quantity.

The sharp constants in systolic inequalities are known only for surfaces
($\TT^2$, $\RP^2$, $\mathrm{Klein\,bottle}$) and are open in higher
dimensions.  The interplay with lattice theory, Hermite constants, and
the geometry of optimal sphere packings suggests deep connections with
number theory and algebraic geometry.

\begin{thebibliography}{99}

\bibitem{Gromov1983}
M.~Gromov,
\textit{Filling Riemannian manifolds},
J. Differential Geom. \textbf{18} (1983), no.~1, 1--147.

\bibitem{Pu1952}
P.~M. Pu,
\textit{Some inequalities in certain nonorientable Riemannian manifolds},
Pacific J. Math. \textbf{2} (1952), 55--71.

\bibitem{Katz1983}
M.~Katz,
\textit{The filling radius of two-point homogeneous spaces},
J. Differential Geom. \textbf{18} (1983), no.~3, 505--511.

\bibitem{Gromov1981}
M.~Gromov,
\textit{Structures métriques pour les variétés riemanniennes},
Textes Mathématiques \textbf{1}, CEDIC, Paris, 1981
(edited by J.\ Lafontaine and P.\ Pansu).

\bibitem{Gromov2003waist}
M.~Gromov,
\textit{Isoperimetry of waists and concentration of maps},
Geom. Funct. Anal. \textbf{13} (2003), 178--215.

\bibitem{Rotman2006}
R.~Rotman,
\textit{The length of a shortest geodesic net on a closed Riemannian manifold},
Topology \textbf{46} (2007), 343--356.

\bibitem{Sabourau2004}
S.~Sabourau,
\textit{Systoles des surfaces plates singulières de genre deux},
Math. Z. \textbf{247} (2004), 693--709.

\bibitem{BergerPansu2012}
M.~Berger,
\textit{A Panoramic View of Riemannian Geometry},
Springer, Berlin, 2003.

\bibitem{CrootLev}
M.~Katz and S.~Sabourau,
\textit{Hyperelliptic surfaces are Loewner},
Proc. Amer. Math. Soc. \textbf{134} (2006), 1189--1195.

\bibitem{KatzBook}
M.~G.~Katz,
\textit{Systolic Geometry and Topology},
Mathematical Surveys and Monographs \textbf{137},
American Mathematical Society, 2007.

\bibitem{BabenkoBrunnbauer}
I.~Babenko,
\textit{Forte souplesse intersystolique de variétés fermées et de polyèdres},
Ann. Inst. Fourier \textbf{52} (2002), 1259--1284.

\end{thebibliography}

\end{document}
